% GENERATED FILE. Edit canonical lessons, then run npm run generate:books.
% canonical-source-hash: fnv1a64:644edb8e53b256e2
% canonical-lessons: UR-C32-pehla, UR-C32-dusra-tisra, UR-C32-chautha, UR-C32-practice, UR-R32-close

\chapter{Which One: the Ordinals}
\label{ch:ur-ordinals}
\hlchaptermodality{32}

\begin{chapteropening}
\textbf{By the end of this chapter:} \emph{I can say first, second, third and fourth, and build the ordinal of any number I can count to.}

Three irregular forms at the bottom of the series and one ending above them, with the break at first set against the same break in English.
\end{chapteropening}

\section[pehlā]{\ur{پہلا} — first}

\label{lesson:UR-C32-pehla}



\begin{quote}
Read \textbf{\ur{۱}}, \textbf{\ur{۵}} and \textbf{\ur{۱۰}}, then say \textbf{ek}.
\end{quote}

\subsection*{You'll want to know first — one word}
\begin{quote}
\textbf{\ur{پہلا}} — \emph{pehlā} — \textbf{first}
\end{quote}

Counting says \emph{how many}. Ordering says \emph{which one} — and a floor number, a date on a notice, a queue, a chapter heading and \textquotedblleft{}the first bus\textquotedblright{} all need the second. Twelve numbers do not give you one of them.

\begin{cousinweb}
\textbf{\ur{پہلا}} does NOT come from \textbf{\ur{ایک}}. It descends through Prakrit from Sanskrit \emph{prathama}, \textquotedblleft{}first, foremost,\textquotedblright{} a separate word.

English does the same thing: \textbf{first} is not built from \emph{one}, and \textbf{second} is not built from \emph{two}, while \textbf{third}, \textbf{fourth} and \textbf{fifth} fall into line behind their numbers. Urdu breaks in the same place, for the same reason — the words people say most often wear away from their families first.
\end{cousinweb}

\subsection*{Your turn}
\begin{itemize}
\raggedright
  \item \emph{Say it:} \textbf{pehlā} — first
  \item \emph{Contrast:} \textbf{ek} counts, \textbf{pehlā} orders
  \item \emph{Compare:} English \textbf{one}/\textbf{first} breaks in the same place
\end{itemize}

\begin{tcolorbox}[breakable,colback=teal!4,colframe=teal!35!black,title={Before you move on}]
Sources: \href{https://en.wiktionary.org/wiki/\%D9\%BE\%DB\%81\%D9\%84\%D8\%A7}{Wiktionary: \ur{پہلا}}; \href{https://www.omniglot.com/language/numbers/urdu.htm}{Omniglot, Numbers in Urdu}.
\end{tcolorbox}
\section[dūsrā tīsrā]{\ur{دوسرا} and \ur{تیسرا} — second, third, and the ending}

\label{lesson:UR-C32-dusra-tisra}



\begin{quote}
Say \textbf{pehlā} and why it is not built from \textbf{ek}. Then say \textbf{do} and \textbf{tīn}.
\end{quote}

\subsection*{You'll want to know first — two words and an ending}
\begin{quote}
\textbf{\ur{دوسرا}} — \emph{dūsrā} — \textbf{second}
\end{quote}

\begin{quote}
\textbf{\ur{تیسرا}} — \emph{tīsrā} — \textbf{third}
\end{quote}

Look at what came back. \textbf{\ur{دو}} is inside \emph{dūsrā}; \textbf{\ur{تین}} is inside \emph{tīsrā}. The number returns, and an ending \emph{-rā} is added to it.

So the series breaks at \emph{first} and mends at \emph{second} — one word out of line, then a pattern. English breaks one word later and mends at \emph{third}; the shape is the same and only the seam moved.

\subsection*{Your turn}
\begin{itemize}
\raggedright
  \item \emph{Say it:} \textbf{dūsrā} — second; \textbf{tīsrā} — third
  \item \emph{Split:} \emph{do} + \emph{-srā}, \emph{tīn} + \emph{-srā}
  \item \emph{Run through:} \textbf{pehlā, dūsrā, tīsrā}
\end{itemize}

\begin{tcolorbox}[breakable,colback=teal!4,colframe=teal!35!black,title={Before you move on}]
Sources: \href{https://en.wiktionary.org/wiki/\%D8\%AF\%D9\%88\%D8\%B3\%D8\%B1\%D8\%A7}{Wiktionary: \ur{دوسرا}}; \href{https://www.omniglot.com/language/numbers/urdu.htm}{Omniglot, Numbers in Urdu}.
\end{tcolorbox}
\section[chauthā]{\ur{چوتھا} — fourth, and the ending that runs all the way up}

\label{lesson:UR-C32-chautha}



\begin{quote}
Say \textbf{dūsrā} and \textbf{tīsrā}, and the number inside each.
\end{quote}

\subsection*{You'll want to know first — one word, then the rule above it}
\begin{quote}
\textbf{\ur{چوتھا}} — \emph{chauthā} — \textbf{fourth}
\end{quote}

\textbf{\ur{چار}} is still audible in it, bent at the vowel.

Above four the language settles. From five upward Urdu adds \emph{-vāṅ} to the plain number and stops bending: \emph{pānchvāṅ}, \textquotedblleft{}fifth\textquotedblright{}; \emph{dasvāṅ}, \textquotedblleft{}tenth\textquotedblright{}; \emph{bīsvāṅ}, \textquotedblleft{}twentieth\textquotedblright{}.

So the whole series is three irregular forms at the bottom — \emph{pehlā}, \emph{dūsrā}, \emph{tīsrā}, with \emph{chauthā} half a step out — and one ending for everything above them. Twelve numbers and one ending is twelve ordinals.

\subsection*{Your turn}
\begin{itemize}
\raggedright
  \item \emph{Say it:} \textbf{chauthā} — fourth
  \item \emph{Build:} \emph{pānch} + \emph{-vāṅ} $\to$ \textbf{pānchvāṅ}, fifth
  \item \emph{Build:} \emph{das} + \emph{-vāṅ} $\to$ \textbf{dasvāṅ}, tenth
\end{itemize}

\begin{tcolorbox}[breakable,colback=teal!4,colframe=teal!35!black,title={Before you move on}]
Sources: \href{https://en.wiktionary.org/wiki/\%DA\%86\%D9\%88\%D8\%AA\%DA\%BE\%D8\%A7}{Wiktionary: \ur{چوتھا}}; \href{https://www.omniglot.com/language/numbers/urdu.htm}{Omniglot, Numbers in Urdu}.
\end{tcolorbox}
\section[Practice]{Payoff — which one}

\label{lesson:UR-C32-practice}



\begin{quote}
Count aloud to ten, read \textbf{\ur{۷}}, \textbf{\ur{۲۰}} and \textbf{\ur{۱۰۰}}, and say which way the digits ran.
\end{quote}

\subsection*{Your turn}
\begin{enumerate}
\raggedright
  \item Say the first four in order, then out of order.
  \item \textbf{Build}, do not recall: fifth, seventh, tenth, twentieth, hundredth.
  \item Say which of the series is not built from its number at all, and say what English does in the same place.
  \item Use one: name the floor a room is on, and the day of a month a notice gives.
  \item Say what the ordinals still cannot reach, and why — the same nine irregular teens that the cardinals could not reach either.
\end{enumerate}

\begin{tcolorbox}[breakable,colback=teal!4,colframe=teal!35!black,title={Before you move on}]
Source: \href{https://www.omniglot.com/language/numbers/urdu.htm}{Omniglot, Numbers in Urdu}.
\end{tcolorbox}
\section[Practice]{Close the numeral slice}

\label{lesson:UR-R32-close}



\begin{quote}
Say the first four ordinals, then count one to five.
\end{quote}

\subsection*{Your turn}
Everything this slice added, with nothing in front of you.

\begin{enumerate}
\raggedright
  \item Count aloud one to ten, then say \textbf{bīs} and \textbf{sau}.
  \item Write the ten figures shuffled, checking \textbf{\ur{۴}}, \textbf{\ur{۵}} and \textbf{\ur{۶}} against the models rather than against memory.
  \item Say the first four ordinals, then build \textbf{fifth}, \textbf{tenth} and \textbf{twentieth}.
  \item Say the two things this slice deliberately did NOT give you, and why: the nine irregular teens, and the ordinals built on them.
\end{enumerate}

Then say what the whole slice cost the hand: \textbf{no new letter at all.} Every sign in every one of the twelve number words was already on the page; only the ten figures were new. And say which road all twelve came by.

\begin{tcolorbox}[breakable,colback=teal!4,colframe=teal!35!black,title={Before you move on}]
Source: \href{https://www.omniglot.com/language/numbers/urdu.htm}{Omniglot, Numbers in Urdu}.
\end{tcolorbox}
